\subsection{评估闭环：分数的用途是做决策，不是装饰榜单}

slides 49--52 先回答评估为何存在：识别一次改动是否有效、在候选模型中做选择、判断系统是否可以进入生产。close-ended evaluation 通过有限答案空间实现自动验证，MMLU 是典型例子；但 prompt 模板、选项顺序、解析器和 train-test contamination 都会影响结果。课程特别提醒，开发期最危险的不是“测试集被模型见过”这一标签本身，而是团队误以为分数变化来自能力，实际上来自评估协议变化。

\lecturefigure{slides-images/slide-049.png}{官方 slide 49：评估服务于改进识别、模型选择与生产决策}
\lecturefigure{slides-images/slide-050.png}{官方 slide 50：MMLU 式 close-ended 自动验证}
\lecturefigure{slides-images/slide-051.png}{官方 slide 51：prompt 敏感性与评估不一致}
\lecturefigure{slides-images/slide-052.png}{官方 slide 52：提示敏感性与 train-test contamination}

\begin{knowledgebox}{读图：先问“这个分数用于哪项决策”}
同一个 benchmark 可以用于快速回归、研究比较或上线门槛，但三种用途要求不同。快速回归重视稳定、便宜和可定位；研究比较需要公开协议与置信区间；上线门槛还要覆盖真实任务分布、延迟、成本和安全。slide 49 的三个动词 Identify、Select、Decide 正是在提醒评估不是单一排行榜。若评估不能改变研发选择，它只是观测报表；若一个分数被用于高风险上线决策，就必须补充失败分布和人工复核。
\end{knowledgebox}

\begin{warningbox}{Contamination 不只有“训练集抄答案”一种形式}
模型可能见过题目、答案、题目改写或 benchmark 讨论；开发者也可能通过反复调 prompt 对测试集过拟合。反过来，发现疑似污染也不代表 benchmark 完全无用：它仍可用于回归，但不能再被解释为纯粹的 unseen generalization。\textbf{老师强调：}课堂说 contamination 对日常开发“没有那么重要”，语境是只要协议稳定，它仍能帮助比较版本；这不等于在科学报告中可以忽略泄漏。
\end{warningbox}

评估协议应被视为版本化代码。模型版本、prompt、采样温度、解码上限、工具版本、判分器和数据快照都要一并记录。Agent 任务还要保存环境状态和执行轨迹，否则同一模型在不同 API 响应或权限配置下产生的差异会被错误归因给训练。close-ended 指标提供速度，但它只能承担可自动判定的那一部分目标。

